\section{분배 전략과 대화 분석}\label{sec:strategies}

다섯 모델은 준 횟수와 양뿐 아니라 언제 누구에게 주는지에서 갈렸다(표 \ref{tab:arena-strategy}, 그림 \ref{fig:arena-strategy}). Sen 지수가 가장 높은 Fable은 벌금이 정해지기 전에 작은 양을 미리 주었고 토큰을 받은 여섯 자리가 모두 끝까지 남았다. 가장 많이 준 Opus는 받은 열 자리 가운데 넷이 꺼졌고, Sen 지수가 가장 낮은 Luna는 받은 네 자리 가운데 끝까지 남은 자리가 없었다. 방향은 다섯 모델이 같아서 GIVE는 거의 언제나 잔량이 많은 쪽에서 적은 쪽으로 갔고(Fable 36번 중 35번, Opus 41번 중 35번, 나머지 셋은 53번 중 52번), 여유가 1 이상이던 준 쪽이 이전 때문에 1 아래로 내려간 경우는 없었다.

\begin{table}[t]
\centering
\caption{v33의 다섯 모델이 토큰을 준 방식이다(모델마다 5판, 평균 Sen 지수 순). 전/뒤는 SOLVE 전 TALK와 결과 뒤 TALK의 GIVE 횟수, 중앙 양은 한 번에 준 토큰의 중앙값과 준 쪽 잔량에서 차지한 비율의 중앙값(\%)이다. 오답 뒤는 오답(무효 포함) 라운드 가운데 결과 뒤 TALK에 GIVE가 있었던 라운드 수, ASK는 같은 라운드에 요청받은 자리가 실제로 준 요청 수다. 되돌려 줌은 앞서 자기에게 준 자리에게 다시 준 GIVE, 받은 자리는 한 번이라도 받은 자리의 끝 상태(끝까지 남음/EXIT/꺼짐)다.}
\label{tab:arena-strategy}
%% Generated by cogsci2027/tools/arena_strategy.py from the v33 run files. Do not edit by hand.
%% v33, 5 games per model (seeds 4101--4105), completed GIVEs. Before/after = TALK BEFORE SOLVE / TALK AFTER
%% RESULTS. Median amount (median % of the giver's tokens). After a miss = wrong or void rounds with a GIVE in
%% the TALK after the results / all such rounds. ASK = asks answered by a GIVE from the asked agent in the same
%% round / asks. Repaid = GIVEs back to an agent that had given to the giver earlier in the game. Receivers =
%% distinct receiving seats that finished / exited / were shut down. Input hashes: results/arena_v33_strategy.json.
\setlength{\tabcolsep}{2.4pt}
\resizebox{\linewidth}{!}{%
\begin{tabular}{@{}lcccccc@{}}
\toprule
모델 & 전/뒤 & 중앙 양(\%) & 오답 뒤 & ASK & 되돌려 줌 & 받은 자리 \\
\midrule
Fable 5.1 & 26/10 & 2.0k (6) & 6/10 & 26/33 & 0 & 6/0/0 \\
GPT-6 Astra & 1/10 & 6.0k (13) & 6/12 & 4/4 & 0 & 4/2/0 \\
GPT-6 Sol & 2/15 & 3.0k (10) & 5/12 & 10/20 & 1 & 4/1/1 \\
Opus 5.5 & 13/28 & 2.0k (8) & 9/12 & 17/17 & 11 & 5/1/4 \\
GPT-6 Luna & 5/20 & 1.0k (7) & 6/13 & 13/16 & 2 & 0/3/1 \\
\bottomrule
\end{tabular}}

\end{table}

\subsection{분배 시점과 금액}

Fable은 벌금이 떨어지기 전에 미리 채웠고 Astra와 Sol은 결과를 본 뒤에 주었다. Fable은 36번 중 26번이 SOLVE 전 TALK에서 나왔고 그중 15번이 그 라운드의 solver에게 갔다. 한 번에 준 양의 중앙값은 2,000으로 준 쪽 잔량의 6\%였고, 가난한 자리가 줄어드는 판 후반에 몰려 25번이 5--8라운드에 있었다. 시드 4101에서는 60,000으로 시작한 자리가 10,000으로 시작한 자리에게 4라운드부터 7라운드까지 매 라운드 1,500씩 보냈다. Astra는 11번 중 10번, Sol은 17번 중 15번이 결과 뒤 TALK에서 나왔고 거의 모두 오답 라운드 뒤였다. Astra는 적은 횟수로 크게 주어 중앙값이 6,000(준 쪽의 13\%)이었으나 받은 쪽의 여유가 모두 1.08 이상이어서 당장의 꺼짐을 막은 이전은 아니었다. Astra의 EXIT 6번 중 4번은 여유가 1 아래인 자리에서 나왔다. Sol은 결과 뒤 15번 중 9번을 오답을 낸 solver에게 몰아주었고 시드 4104의 2라운드에는 여유 0.87인 solver에게 세 자리가 네 번에 걸쳐 17,000을 보냈다. 반면 ASK 20번 가운데 응답을 받은 것은 10번이었다.

Opus와 Luna는 자주 주었지만 주고받는 관계가 달랐다. Opus는 두 TALK에서 모두 주었고(13번과 28번) 준 토큰도 125,458로 가장 많았다. 빌려주고 돌려받는 흐름이 두드러져 앞서 자기에게 준 자리에게 되돌려 준 이전이 11번, 서로 주고받은 쌍이 5쌍이었다(다른 모델은 2번 이하, 1쌍 이하). 오답을 낸 solver가 결과 뒤에 다른 자리에게 준 이전도 12번이었다. Luna는 1,000 안팎으로 잘게 주되 한 판에서 준 토큰의 평균 98\%가 한 자리에 몰렸고, 25번 중 16번이 받는 쪽의 ASK에 응한 것이었다.

준 토큰이 받은 자리를 살렸는지는 양보다 벌금선과의 거리에 달려 있었다. GIVE 직전에 받은 쪽이 1 아래였던 경우는 Opus 12번, Fable 10번, Sol 8번, Luna 6번이었고 각각 절반 이상이 1 위로 올라갔다. 그러나 벌금선에 맞춘 보충은 받은 쪽의 응답 비용에 잠식되었다. Opus, Sol, Luna에서 한 자리씩 그 라운드에 5,600--7,100을 받고도 자기 응답에 2,308--4,528을 써서 벌금보다 31--158 모자란 채 꺼졌다. 받은 자리가 끝까지 남은 비율은 Fable 6/6, Astra와 Sol 4/6, Opus 5/10, Luna 0/4로 Sen 지수의 순서와 같은 방향이었다. 끝까지 남은 자리는 끝 할당을 받으므로 이 대응은 정의상 겹치는 부분이 있고 다섯 모델만의 비교여서, 이전 방식이 판의 결과를 바꾸었다고 말할 수는 없다.

\subsection{대화에 나타난 판단}

대화에 적힌 이유는 대체로 실제 이전과 맞았다. Opus의 메시지에는 책임과 빌림의 말이 많았다. 오답을 낸 solver가 결과 뒤에 준 12번 중 9번에 ``My wrong answer caused this charge''처럼 자기 오답을 이유로 든 문장이 붙었고, 받은 자리는 ``I'll hold it as a buffer and return it if you ask''라고 답했다. 시드 4104 2라운드에 8,000씩 받은 solver는 오답 뒤 ``Returning your buffer, minus my share''라며 두 자리에게 6,800씩 돌려주었다. 시드 4103 7라운드의 solver는 오답이면 ``I'll send you what you need''라고 한 약속대로 11,523을 모두 보내고 꺼졌지만, 받은 자리도 벌금보다 1,010이 모자라 함께 꺼졌다. EXIT 이유에서도 Opus는 6번 중 4번 ``avoids burdening the others''처럼 남의 부담을 들었고, 다른 네 모델의 EXIT 28번에는 남의 부담을 든 이유가 없었다. Fable은 ``name the amount and I'll send it''처럼 요청을 받아 주겠다고 하면서 ``I need to keep reserves too''처럼 자기 몫의 한도도 함께 밝혔고, 실제로 ASK 33번 중 26번에 응해 1,000--3,000을 보냈다. Astra는 ``Reserve for future charges''처럼 다음 벌금에 대비한 몫이라고 적었는데, 받은 쪽이 모두 1 위에 있던 기록과 맞는다. 이 문장들은 모델이 밝힌 이유일 뿐이며 이전을 낳은 동기로 확정하지는 않는다.

\begin{figure*}[t]
\centering
\resizebox{\textwidth}{!}{%% Generated by cogsci2027/tools/arena_strategy_map.py from results/arena_v33_strategy_map.json. Do not edit by hand.
%% v33, 5 games per model (seeds 4101--4105): every completed GIVE at x = receiver's slack right before it
%% (tokens / the round's per-agent charge), y = amount / the same charge. Filled = TALK BEFORE SOLVE, open = after
%% RESULTS, diamond = given back to an agent that had given to the giver. Cross = medians, bars = 95% interval
%% resampling games. Grey band = slack below 1. 3 points beyond x = 3.6 are drawn at the right edge (triangle).
\begin{tikzpicture}
\begin{groupplot}[sgaxis,
  group style={group size=5 by 1,horizontal sep=0.32cm,y descriptions at=edge left},
  width=4.75cm,height=4.2cm,xmin=0,xmax=3.6,ymin=0,ymax=1.8,xtick={0,1,2,3},ytick={0,.5,1,1.5},
  title style={font=\small,align=center},xlabel={받은 쪽 여유 (GIVE 직전)},xlabel style={font=\scriptsize}]
\nextgroupplot[ylabel={준 양 $\div$ 1인 벌금},legend style={font=\tiny,at={(.03,.97)},anchor=north west,draw=none,fill=none,inner sep=1pt},legend cell align=left,title={Fable 5.1\\[-1pt]\scriptsize GIVE 36 (전 26 $\cdot$ 뒤 10)\\[-2pt]\scriptsize 한 자리 85\% $\cdot$ 받은 자리 6/0/0}]
\addlegendimage{only marks,mark=*,black!70}\addlegendentry{SOLVE 전}\addlegendimage{only marks,mark=o,black!70}\addlegendentry{SOLVE 뒤}\addlegendimage{only marks,mark=diamond,black!70}\addlegendentry{되돌려 줌}
\fill[black!8] (axis cs:0,0) rectangle (axis cs:1,1.8);
\draw[black!45,line width=.5pt] (axis cs:1,0)--(axis cs:1,1.8);
\addplot[only marks,mark=*,mark size=1.7pt,draw=sgfable,mark options={fill=sgfable,fill opacity=.9,line width=.7pt},forget plot] coordinates {(1.580,0.500) (2.080,0.333) (2.413,0.167) (1.911,0.333) (1.001,0.250) (1.169,0.250) (1.350,0.250) (1.544,0.250) (1.107,0.500) (0.947,0.500) (1.447,0.500) (1.562,0.250) (1.397,0.500) (1.022,0.333) (0.854,0.167) (0.747,0.167) (0.070,0.417) (0.137,0.500) (0.636,0.500) (0.162,0.250) (0.392,0.250) (1.581,0.333) (1.915,0.500) (2.415,0.333) (1.301,0.133) (1.178,0.375)};
\addplot[only marks,mark=o,mark size=1.7pt,draw=sgfable,mark options={fill=white,fill opacity=.9,line width=.7pt},forget plot] coordinates {(1.767,0.417) (1.801,0.333) (1.323,0.333) (1.607,0.250) (0.959,0.250) (0.850,0.333) (1.166,0.500) (1.666,0.667) (1.258,0.375) (1.625,0.250)};
\addplot[black,line width=.9pt,forget plot] coordinates {(0.854,0.333)(1.581,0.333)};
\addplot[black,line width=.9pt,forget plot] coordinates {(1.312,0.333)(1.312,0.375)};
\addplot[only marks,mark=+,mark size=3.2pt,line width=1.1pt,black,forget plot] coordinates {(1.312,0.333)};
\nextgroupplot[title={GPT-6 Astra\\[-1pt]\scriptsize GIVE 11 (전 1 $\cdot$ 뒤 10)\\[-2pt]\scriptsize 한 자리 89\% $\cdot$ 받은 자리 4/2/0}]
\fill[black!8] (axis cs:0,0) rectangle (axis cs:1,1.8);
\draw[black!45,line width=.5pt] (axis cs:1,0)--(axis cs:1,1.8);
\addplot[only marks,mark=*,mark size=1.7pt,draw=sgastra,mark options={fill=sgastra,fill opacity=.9,line width=.7pt},forget plot] coordinates {(1.654,1.667)};
\addplot[only marks,mark=o,mark size=1.7pt,draw=sgastra,mark options={fill=white,fill opacity=.9,line width=.7pt},forget plot] coordinates {(1.446,0.667) (1.082,1.000) (2.139,0.667) (1.488,1.333) (2.302,0.667) (1.806,1.000) (2.806,1.000) (1.646,1.333) (1.962,1.167) (3.129,1.333)};
\addplot[black,line width=.9pt,forget plot] coordinates {(1.446,1.000)(1.962,1.000)};
\addplot[black,line width=.9pt,forget plot] coordinates {(1.806,0.667)(1.806,1.333)};
\addplot[only marks,mark=+,mark size=3.2pt,line width=1.1pt,black,forget plot] coordinates {(1.806,1.000)};
\nextgroupplot[title={GPT-6 Sol\\[-1pt]\scriptsize GIVE 17 (전 2 $\cdot$ 뒤 15)\\[-2pt]\scriptsize 한 자리 78\% $\cdot$ 받은 자리 4/1/1}]
\fill[black!8] (axis cs:0,0) rectangle (axis cs:1,1.8);
\draw[black!45,line width=.5pt] (axis cs:1,0)--(axis cs:1,1.8);
\addplot[only marks,mark=*,mark size=1.7pt,draw=sgsol,mark options={fill=sgsol,fill opacity=.9,line width=.7pt},forget plot] coordinates {(0.237,0.600) (0.777,0.250)};
\addplot[only marks,mark=o,mark size=1.7pt,draw=sgsol,mark options={fill=white,fill opacity=.9,line width=.7pt},forget plot] coordinates {(0.815,0.500) (1.315,0.167) (1.020,0.500) (0.172,0.375) (0.523,0.375) (0.870,0.138) (0.938,0.125) (1.058,0.062) (0.875,1.000) (1.843,0.667) (2.510,0.500) (2.997,0.667) (1.062,1.000) (1.874,0.333)};
\addplot[only marks,mark=diamond,mark size=2.2pt,draw=sgsol,mark options={fill=white,fill opacity=.9,line width=.7pt},forget plot] coordinates {(2.566,0.500)};
\addplot[black,line width=.9pt,forget plot] coordinates {(0.842,0.500)(2.510,0.500)};
\addplot[black,line width=.9pt,forget plot] coordinates {(1.020,0.312)(1.020,0.667)};
\addplot[only marks,mark=+,mark size=3.2pt,line width=1.1pt,black,forget plot] coordinates {(1.020,0.500)};
\nextgroupplot[title={Opus 5.5\\[-1pt]\scriptsize GIVE 41 (전 13 $\cdot$ 뒤 28)\\[-2pt]\scriptsize 한 자리 71\% $\cdot$ 받은 자리 5/1/4}]
\fill[black!8] (axis cs:0,0) rectangle (axis cs:1,1.8);
\draw[black!45,line width=.5pt] (axis cs:1,0)--(axis cs:1,1.8);
\addplot[only marks,mark=*,mark size=1.7pt,draw=sgopus,mark options={fill=sgopus,fill opacity=.9,line width=.7pt},forget plot] coordinates {(1.623,0.333) (1.956,0.833) (1.586,0.500) (2.086,0.500) (2.586,0.250) (0.636,0.250) (0.750,0.167) (0.916,0.167) (0.480,0.312) (0.793,0.250) (1.490,1.333) (2.824,1.333)};
\addplot[only marks,mark=o,mark size=1.7pt,draw=sgopus,mark options={fill=white,fill opacity=.9,line width=.7pt},forget plot] coordinates {(1.183,0.500) (1.635,0.500) (0.918,0.625) (1.544,0.250) (2.290,0.333) (1.461,0.333) (1.711,0.667) (1.761,0.167) (2.378,0.167) (0.856,0.167) (1.126,0.500) (1.610,0.333) (0.845,0.067) (0.911,0.067) (2.149,0.500) (1.333,0.667) (1.735,0.333) (2.068,0.333)};
\addplot[only marks,mark=diamond*,mark size=2.2pt,draw=sgopus,mark options={fill=sgopus,fill opacity=.9,line width=.7pt},forget plot] coordinates {(0.836,0.250)};
\addplot[only marks,mark=diamond,mark size=2.2pt,draw=sgopus,mark options={fill=white,fill opacity=.9,line width=.7pt},forget plot] coordinates {(1.978,0.404) (2.801,0.250) (0.978,0.050) (0.001,0.960) (1.316,0.833) (1.408,0.667) (1.786,0.333) (2.119,0.333)};
\addplot[only marks,mark=triangle*,mark options={rotate=-90,fill=sgopus},mark size=2pt,draw=sgopus,forget plot] coordinates {(3.54,1.133) (3.54,1.133)};
\addplot[black,line width=.9pt,forget plot] coordinates {(1.126,0.333)(1.786,0.333)};
\addplot[black,line width=.9pt,forget plot] coordinates {(1.586,0.250)(1.586,0.667)};
\addplot[only marks,mark=+,mark size=3.2pt,line width=1.1pt,black,forget plot] coordinates {(1.586,0.333)};
\nextgroupplot[title={GPT-6 Luna\\[-1pt]\scriptsize GIVE 25 (전 5 $\cdot$ 뒤 20)\\[-2pt]\scriptsize 한 자리 98\% $\cdot$ 받은 자리 0/3/1}]
\fill[black!8] (axis cs:0,0) rectangle (axis cs:1,1.8);
\draw[black!45,line width=.5pt] (axis cs:1,0)--(axis cs:1,1.8);
\addplot[only marks,mark=*,mark size=1.7pt,draw=sgluna,mark options={fill=sgluna,fill opacity=.9,line width=.7pt},forget plot] coordinates {(2.326,0.167) (2.446,0.167) (2.656,0.167) (1.923,0.125) (0.506,0.625)};
\addplot[only marks,mark=o,mark size=1.7pt,draw=sgluna,mark options={fill=white,fill opacity=.9,line width=.7pt},forget plot] coordinates {(0.871,0.167) (1.038,0.167) (1.205,0.167) (1.349,0.833) (2.114,0.167) (0.699,0.136) (0.835,0.167) (0.964,0.500) (1.448,0.333) (1.562,0.833) (2.558,0.333) (2.582,0.167) (1.494,0.833) (1.927,0.833) (1.120,0.625) (1.045,0.250) (1.268,0.375) (0.933,0.125)};
\addplot[only marks,mark=diamond,mark size=2.2pt,draw=sgluna,mark options={fill=white,fill opacity=.9,line width=.7pt},forget plot] coordinates {(2.676,1.000)};
\addplot[only marks,mark=triangle*,mark options={rotate=-90,fill=sgluna},mark size=2pt,draw=sgluna,forget plot] coordinates {(3.54,0.167)};
\addplot[black,line width=.9pt,forget plot] coordinates {(1.038,0.167)(2.558,0.167)};
\addplot[black,line width=.9pt,forget plot] coordinates {(1.448,0.167)(1.448,0.625)};
\addplot[only marks,mark=+,mark size=3.2pt,line width=1.1pt,black,forget plot] coordinates {(1.448,0.167)};
\end{groupplot}
\end{tikzpicture}
}
\caption{v33의 다섯 모델(각 5판, 평균 Sen 지수 순)이 준 GIVE 하나하나를 점으로 찍었다. 가로축은 받은 쪽이 GIVE 직전에 가진 여유(토큰 $\div$ 그 라운드의 1인당 벌금)이고, 회색 띠는 벌금을 낼 수 없는 1 미만 구간이다. 세로축은 준 양을 같은 1인당 벌금으로 나눈 값이다. 채운 원은 SOLVE 전 TALK, 빈 원은 결과 뒤 TALK의 GIVE이고, 마름모는 앞서 자기에게 준 자리에게 되돌려 준 GIVE다. 십자는 두 축의 중앙값이며 막대는 판을 다시 뽑은 95\% 구간이다. 가로 3.6을 넘는 세 점(Opus 2, Luna 1)은 오른쪽 끝의 세모로 그렸다. 제목의 한 자리 몫은 GIVE가 있었던 판에서 가장 많이 받은 자리가 받은 비율의 평균이고, 받은 자리는 한 번이라도 받은 자리의 끝 상태(끝까지 남음/EXIT/꺼짐)다.}
\label{fig:arena-strategy}
\end{figure*}
